% Business-Paper -- Gerüst der Harwness-Vorlage (Skill latex-report).
% Aufbau nach dem Pyramidenprinzip (Skill business-writing-pyramid):
% Kernaussage zuerst, Einleitung aus Lage, Störung und Frage, darunter die
% tragenden Gründe, jeder mit Belegen. Der Titel ist eine Aussage, kein Thema.
% Alle Texte unten sind kurze Platzhalter: ersetzen, kürzen oder streichen.
% Nichts erfinden -- fehlende Zahlen bleiben als % TODO: stehen.
\documentclass[11pt,a4paper]{scrartcl}

% ── Anpassung: von latex.template gefüllt, Vorgaben in harw-report.sty ──
\newcommand{\harwLanguage}{%%LANGUAGE%%}
\newcommand{\harwAccentHex}{%%ACCENT%%}
\newcommand{\harwWarnHex}{%%WARN%%}
\newcommand{\harwMainFont}{%%MAINFONT%%}
\newcommand{\harwSansFont}{%%SANSFONT%%}
\newcommand{\harwMonoFont}{%%MONOFONT%%}
\usepackage{harw-report}

\title{%%TITLE%%}
\subtitle{%%SUBTITLE%%}
\author{%%AUTHOR%%}
\date{%%DATE%%}

\begin{document}
\maketitle

\begin{merke}[Die Idee in drei Sätzen]
Kernaussage: die Antwort auf die Frage der Leserin in einem Satz. Warum:
die zwei oder drei tragenden Gründe in einem Satz. Was jetzt: die
Entscheidung oder der nächste Schritt, den das Paper vorschlägt.
\end{merke}

\begin{achtung}[Wichtig vorab]
Die eine Einschränkung, die die Leserin kennen muss, bevor sie entscheidet
(etwa eine offene Annahme oder ein vorläufiger Datenstand).
\end{achtung}

\tableofcontents
\newpage

\section{Ausgangslage, Störung, Frage}\label{sec:einleitung}

\begin{description}[leftmargin=2.6cm,style=nextline,font=\sffamily\bfseries\small]
  \item[Lage] Was die Leserin schon weiß und für richtig hält.
  \item[Störung] Was sich geändert hat oder nicht mehr passt.
  \item[Frage] Die eine Frage, die daraus folgt und die dieses Paper beantwortet.
\end{description}

\section{Die Empfehlung}\label{sec:empfehlung}

Die Antwort auf die Frage, ausführlicher als oben: ein Absatz, der die
Kernaussage wiederholt und die drei Gründe ankündigt (Abschnitte
\ref{sec:grund1} bis \ref{sec:grund3}).

\begin{figure}[htbp]
\centering
\begin{tikzpicture}[node distance=0.9cm and 0.6cm]
  \node[harwbox,fill=harwaccent!15,text width=6cm] (kern)
    {\textbf{Kernaussage}\\ein Satz};
  \node[harwbox,fill=harwgrau,below=of kern] (g2) {Grund 2};
  \node[harwbox,fill=harwgrau,left=of g2] (g1) {Grund 1};
  \node[harwbox,fill=harwgrau,right=of g2] (g3) {Grund 3};
  \foreach \g in {g1,g2,g3} \draw[harwpfeil] (kern) -- (\g);
\end{tikzpicture}
\caption{Die Argumentation auf einen Blick (Beispiel)}\label{fig:pyramide}
\end{figure}

\section{Grund 1 als ganzer Satz}\label{sec:grund1}

Die Aussage des Abschnitts steht im Titel. Der Text liefert die Belege,
vom wichtigsten zum unwichtigsten.

\section{Grund 2 als ganzer Satz}\label{sec:grund2}

Belege zu Grund 2. Zahlen in Tabellen stehen rechtsbündig:

\begin{tabularx}{\textwidth}{@{}Lrr@{}}
\toprule
\textbf{Posten} & \textbf{Jahr 1} & \textbf{Jahr 2} \\
\midrule
Posten A & 12\,000 & 9\,500 \\
Posten B & 3\,400 & 3\,400 \\
\midrule
Summe & 15\,400 & 12\,900 \\
\bottomrule
\end{tabularx}

\section{Grund 3 als ganzer Satz}\label{sec:grund3}

Belege zu Grund 3.

\begin{beispiel}
Ein konkretes Beispiel, das den Grund greifbar macht.
\end{beispiel}

\section{Umsetzung und Risiken}\label{sec:umsetzung}

{\small
\begin{longtable}{@{}>{\raggedright}p{3.4cm}>{\raggedright}p{5.4cm}>{\raggedright\arraybackslash}p{6.2cm}@{}}
\toprule
\textbf{Schritt} & \textbf{Risiko} & \textbf{Gegenmaßnahme} \\
\midrule
\endhead
\bottomrule
\endlastfoot
Schritt 1 & Was schiefgehen kann. & Wie das Risiko begrenzt wird. \\
Schritt 2 & Was schiefgehen kann. & Wie das Risiko begrenzt wird. \\
\end{longtable}
}

\section{Stand und Entscheidungsbedarf}\label{sec:stand}

\begin{description}[leftmargin=2.6cm,style=nextline,font=\sffamily\bfseries\small]
  \item[Stand] Was belegt ist.
  \item[Offen] Welche Annahmen noch zu prüfen sind.
  \item[Entscheidung] Was die Leserin bis wann entscheiden soll.
\end{description}

\appendix
\section{Glossar}\label{sec:glossar}

\begin{tabularx}{\textwidth}{@{}lL@{}}
\toprule
\textbf{Begriff} & \textbf{Bedeutung} \\
\midrule
Begriff A & Eine kurze Erklärung in einem Satz. \\
\bottomrule
\end{tabularx}

\begin{thebibliography}{9}
\addcontentsline{toc}{section}{\refname}
\small
\bibitem{quelle1} Nachname, Vorname: \emph{Titel der Quelle}. Verlag, Jahr.
  \url{https://example.org/quelle}
\end{thebibliography}

\end{document}
